\subsection{\texorpdfstring{Dual of the space \(L_{F}^{1}\) (F a separable Banach space)}{Dual of the space L sub F superscript 1 (F a separable Banach space)}}

PROPOSITION 10. --- Let F be a separable Banach space. For every function \(\mathbf{f} \in \overline{\mathcal{L}}_{\mathrm{F}}^{1}\) and every function \(\mathbf{g} \in \mathcal{L}_{\mathrm{F}'_s}^{\infty}\) , the numerical function \(\langle \mathbf{f}, \mathbf{g} \rangle : t \mapsto \langle \mathbf{f}(t), \mathbf{g}(t) \rangle\) is essentially \(\mu\) -integrable, and

\[
\left| \int \langle \mathbf {f}, \mathbf {g} \rangle d \mu \right| \leqslant \overline {{{{N}}}} _ {1} (\mathbf {f}) N _ {\infty} (\mathbf {g}).\tag{3}
\]

For every class \(\dot{\mathbf{g}}\in\mathrm{L}_{\mathrm{F}^{\prime}s}^{\infty}\) , let \(\theta(\dot{\mathbf{g}})\) be the continuous linear form on \(\mathrm{L}_{\mathrm{F}}^{1}\) deduced from the linear form \(\mathbf{f}\mapsto\int\langle\mathbf{f},\mathbf{g}\rangle d\mu\) on \(\overline{\mathcal{L}}_{\mathrm{F}}^{1}\) by passage to the quotient; then \(\theta\) is a linear isometry of \(\mathrm{L}_{\mathrm{F}^{\prime}s}^{\infty}\) onto the strong dual \((\mathrm{L}_{\mathrm{F}}^{1})^{\prime}\) of the Banach space \(\mathrm{L}_{\mathrm{F}}^{1}\) .

For every compact subset K of T and every \(\varepsilon > 0\) , there exists a compact subset \(K'\) of K such that \(\mu(K - K') \leqslant \varepsilon\) and the restriction of f (resp. g) to \(K'\) is a continuous mapping of \(K'\) into F (resp. into \(F'_s\) ); consequently \(\mathbf{g}(K')\) is weakly compact, hence equicontinuous on \(F'\) (TVS, III, §4, No.~2, Th. 1 or IV, §1, Exer. 10). Now, the restriction of the canonical bilinear form on \(F \times F'\) to the product of F and an equicontinuous subset of \(F'\) is continuous for the product topology of the topology of F and \(\sigma(F', F)\) (GT, X, §2, No.~1, Cor. 4 of Prop. 1); it follows that the restriction of \(\langle f, g \rangle\) to \(K'\) is continuous, hence that \(\langle f, g \rangle\) is \(\mu\) -measurable. Moreover,

\[
| \langle \mathbf {f} (t), \mathbf {g} (t) \rangle | \leqslant | \mathbf {f} (t) | \cdot | \mathbf {g} (t) | \leqslant | \mathbf {f} (t) | N _ {\infty} (\mathbf {g})
\]

locally almost everywhere, consequently \(\langle f,g\rangle\) is essentially \(\mu\) -integrable and the inequality (3) holds (Ch. IV, §5, No.~6, Th. 5 and Ch. V, §1, No.~3, Lemma).

It remains to show that \(\theta\) is a surjective isometry. Let u be a continuous linear form on \(L_{F}^{1}\) . The mapping \((h,\mathbf{z})\mapsto u(\widetilde{h}\mathbf{z})\) is a continuous bilinear form on \(L^{1}\times F\) , because

\[
\left| u \left(\widetilde {h} \mathbf {z}\right) \right| \leqslant \| u \| \cdot \mathrm{N} _ {1} (h \mathbf {z}) \leqslant \| u \| \cdot | \mathbf {z} | \cdot \mathrm{N} _ {1} (h).
\]

By Cor. 1 of Th. 1 of No.~5, there exists a mapping g of T into F', belonging to \(L_{F_{s}^{\infty}}\) , such that \(|\mathbf{g}(t)| \leqslant \|u\|\) for all \(t \in T\) and such that \(u(\widetilde{h}\mathbf{z}) = \int \langle h\mathbf{z}, \mathbf{g} \rangle d\mu\) for every function \(h \in L^{1}\) with class \(\widetilde{h}\) in \(L^{1}\) and every \(z \in F\) . In other words, the linear forms u and \(\theta(\dot{\mathbf{g}})\) coincide on the subspace of \(L_{F}^{1}\) generated by the elements of the form \(\widetilde{h}\mathbf{z}\) ( \(\widetilde{h} \in L^{1}\) , \(z \in F\) ). Since this subspace is dense in \(L_{F}^{1}\) (Ch. IV, §3, No.~5, Prop. 10), it follows that \(u = \theta(\dot{\mathbf{g}})\) , which already proves that \(\theta\) is surjective. Moreover, by (3), \(\|\theta(\dot{\mathbf{g}})\| \leqslant N_{\infty}(\mathbf{g}) \leqslant \|u\| = \|\theta(\dot{\mathbf{g}})\|\) , whence \(\|\theta(\dot{\mathbf{g}})\| = N_{\infty}(\mathbf{g})\) , and this concludes the proof.
% semantic-fragments: legacy-input-seam
\relax{}
